\section{Parallel inputs and the scale of a single defect}
\label{sec:parallel}

The operator parent developed several other observables before the Green
branch separated. Their transferable content explains both useful
comparisons and the remaining gap. Localized correlations (OC) and
collision constraints (OA) retain separate manuscripts, registers and
research questions \cite{GNPCorrelations,GNPCollisions}. We do not treat
progress in one observable as an estimate of another.

\subsection{Smooth transverse moments and the Green scale}
For a nonnegative compactly supported \(C^2\) function \(a\), set
\(Q_a=\sum_j a(x_j)y_j^2\), counting every root and multiplicity.
A conjugate pair at \(X\pm\ii d\), where \(a(X)=1\), contributes
\(2d^2\). In reverse Newman time, at simple relevant roots the paired
zero dynamics gives
\begin{equation}\label{eq:smooth-kernel}
 \frac{\dd Q_a}{\dd\tau}=2\sum_{j<k}'K_f(z_j,z_k),\qquad
 K_f=\int_0^1\frac{2a(x_t)e^2-a''(x_t)y_t^2d_x^2}
                         {d_x^2+e^2}\dd t,
\end{equation}
where \(f=a(x)y^2\), \(d_x=x_j-x_k\), \(e=y_j-y_k\) and
\(z_t=z_k+t(z_j-z_k)\). To derive it, differentiate \(f\) along the
root velocities and pair the two interaction rows. The mixed Hessian
terms cancel. In \(|y|\le1\),
\(|K_f|\le\max(2\|a\|_\infty,\|a''\|_\infty)\): no internal gap
appears in this bound. The prime requires complete opposite-pair
summation. The count \eqref{eq:count} makes its far velocity tail
\(O(Z\log U/U)\) on \(|z|\le Z\) after cutoff \(U\ge2Z+2\).
This follows by combining \((z-\zeta)^{-1}+(z+\zeta)^{-1}
=2z/(z^2-\zeta^2)\) before partial summation.

The parent proves, from the canonical unit counts and safe-line
resolvent estimates, that a common \(A\ge1\) gives
\begin{equation}\label{eq:smooth-modulus}
 |Q_{X,R}(\lambda_2)-Q_{X,R}(\lambda_1)|
 \le344A^2\|\chi\|_{(2)}R^{-2}\log^2(2+X)|\lambda_2-\lambda_1|,
\end{equation}
for \(X\ge10\), \(0<R\le1\), \(0\le\lambda_i\le1/2\), where
\(a(x)=\chi((x-X)/R)\), \(0\le\chi\le1\),
\(\chi\in C_c^2([-1,1])\), and
\(\|\chi\|_{(2)}=\sum_{j=0}^2\|\chi^{(j)}\|_\infty\)
\cite{SmoothLocalizationProof,SmoothLocalizationAssessment}.
The proof groups the finitely many roots near the support and compares
the exterior with \(H'/H(x+4\ii)\). The exterior difference has an
absolute inverse-square tail; the unpaired harmonic sum is not bounded
absolutely. Local heat-polynomial splitting isolates the finitely many
multiple events in a fixed compact space-time box. Integrating the same
derivative bound between them and using multiset continuity proves
\eqref{eq:smooth-modulus} through collisions. This argument allows a
derivative jump, but no atom in the first distributional derivative.

In particular an off-axis pair at \(\lambda=0\) yields the conditional
threshold lower bound
\(\Lambda\ge2d^2/[344A^2\|\chi\|_{(2)}\log^2(2+X)]\), with
\(R=1\) and \(\chi(0)=1\). At the threshold every zero is real,
since a nonreal root would persist above it. Apply the modulus between
zero and \(\Lambda\). This lower bound is not a contradiction and is
not an upper bound on the lifetime of every possible defect.

The Green observable has a different local scale. Differentiating its
explicit kernel and substituting \(t=\sinh(au)\) gives, for \(0\le y\le1\),
\[
 \int_\R\partial_yG(u,y)\dd u
 =\int_\R\frac{\cos(ay)}{t^2+\cos^2(ay)}\dd t=\pi.
\]
Since \(G(u,0)=0\),
\begin{equation}\label{eq:Green-mass-depth}
 \int_\R G(u,y)\dd u=\pi y.
\end{equation}
Thus a defect of depth \(d\) has a signal of order \(d/T\) in a
normalized center average over length \(T\). At a newly born double
pair the positive ordinate is of order \(\sqrt\tau\), so the Green
signal has that scale, while the smooth moment has scale \(\tau\).
Collision-stable Lipschitz bounds for the latter do not transfer to
the former.

\subsection{What uniform density and all-order averages do provide}
The parent also obtains, for each fixed \(0<b<1/2\),
\begin{equation}\label{eq:inherited-density}
 N_\lambda(T,v):=\#\{j:T\le x_j\le2T,\ |y_j|>v\}
 \le B_bT\log T\,\ee^{-bv\log T},
\end{equation}
uniformly for \(T\ge8\), \(0\le\lambda\le1/2\), \(0<v<1\)
\cite{SmoothLocalizationProof,SmoothLocalizationAssessment}.
Its imported time-zero input is the density estimate of
Maples--Rodgers, with its constant uniform as the strip boundary
approaches the critical line \cite[Theorem 2.4 and Section 4 (18)]{MaplesRodgers2015}.
Its transport step uses a broad window \(a_T\) and an even convex
\(g\) with \(g(0)=0\). In Newman time the favorable \(a_Tg''\)
term of the chord kernel can be discarded, and conjugate exterior
roots have a favorable vertical sign at horizontal separation at least
two. The remaining cost is
\(D_\chi\|g\|_\infty(\lambda_2-\lambda_1)\log^3T\), independent
of \(g'\) and \(g''\). Take \(g(y)=\cosh(b\log T\,y)-1\), apply
the time-zero tail by layer cake, and absorb
\(T^{b-1}\log^2T\). Markov's inequality gives
\eqref{eq:inherited-density}; small \(v\log T\) is paid by the unit count.

Layer cake now gives the explicit consequence
\[
 \sum_{T\le x_j\le2T}|y_j|^{2p}
 \le B_bT\log T\frac{(2p)!}{(b\log T)^{2p}},\qquad p\ge1.
\]
For local windows the root population is at most \(J\log T\), with
a fixed canonical \(J\). Weighted H\"older, followed by integration
over centers (each root occupies center measure at most \(R\int\chi\)),
therefore yields
\begin{equation}\label{eq:moment-all-orders}
 \frac1T\int_T^{2T}Q_{X,R}^p\dd X
 \le C_bR A_b^p\frac{(2p)!}{(\log T)^p}.
\end{equation}
A fixed enlargement of the dyadic band pays support endpoints. The
constants can be chosen independent of \(p,\lambda,T,R\) in the
declared ranges. This is a useful simultaneous all-order statement;
the factorial cost must remain when choosing a growing order.

None of these bounds removes the exceptional-center problem. For
example \(p=1\) bounds a normalized smooth average at order
\(1/\log T\), whereas one pair contributes only order \(d^2/T\).
Similarly, applying \eqref{eq:inherited-density} to the first absolute
transverse moment gives a Green average allowance of constant order,
too large to resolve \eqref{eq:Green-mass-depth}'s single-defect scale.
The inherited resolution assessment makes this comparison for its
specific density-based allowance \cite{OGResolution}. It is a scoped
failure of that allowance, not a lower bound for the actual Green value
or an impossibility theorem for all statistics.

\subsection{A field-based comparison: Gaussian Laguerre defects}
Another parent observable is
\[
 A_n(\lambda,X,v)=\int_\R w_v(x-X)[(f_n')^2-f_nf_n'']\dd x,
 \quad f_n=\partial_x^nH_\lambda,\quad
 w_v(t)=\frac{\ee^{-t^2/(2v)}}{\sqrt{2\pi v}},\quad v>0.
\]
The complete source gives the exact signed form
\begin{equation}\label{eq:Gaussian-source}
 A_n=\frac18\int_{\R^2}(-pq)^n(p-q)^2W_\lambda(p)W_\lambda(q)
             \ee^{\ii X(p+q)-v(p+q)^2/2}\dd p\dd q.
\end{equation}
Symmetrize \((f_n')^2-f_nf_n''\) under \(p\leftrightarrow q\) to get
\((p-q)^2/8\); Gaussian Fourier integration gives the last factor.
Complete source moments justify Fubini. Positivity of the theta factors
does not remove the signed phase or the factor \((-pq)^n\).

Direct differentiation gives
\begin{equation}\label{eq:Gaussian-hierarchy}
 \partial_\lambda A_n=-\partial_{XX}A_n+2A_{n+1},\qquad
 \partial_v A_n=\tfrac12\partial_{XX}A_n.
\end{equation}
If \(M_j(L)=\int|u|^j\ee^{Lu^2}\Phi(u)\dd u\), then
\(|pq|^n\le n!\epsilon^{-n}\ee^{\epsilon(p^2+q^2)/2}\) proves
\( |A_n|\le n!\epsilon^{-n}M_0(L+\epsilon/2)M_2(L+\epsilon/2)/4\)
for \(\lambda\le L\). Consequently the characteristic
\(v=v_0+2h>0\) has the convergent expansion
\[
 A_n(\lambda_0+h,X,v_0+2h)
 =\sum_{k\ge0}\frac{2^kh^k}{k!}A_{n+k}(\lambda_0,X,v_0).
\]
It can also be proved directly by expanding \(\ee^{-2hpq}\) in
\eqref{eq:Gaussian-source}, with complete moment domination. The parent
records explicit joint order, arithmetic and radial remainders and the
assessment of those fixed statements \cite{GaussianLaguerreProof,LocalizedMetricsAssessment}.

At a real event of finite multiplicity \(m\), let
\(a_m=H_{\lambda_*}^{(m)}(x_*)\ne0\). Differentiating the local heat
polynomial \(m-2\) times leaves
\(g(t,y)=a_m(y^2/2+t)+\text{remainder}\), for
\(g=\partial_x^{m-2}H_{\lambda_*-t}(x_*+y)\).
Uniform real-line bounds on \(g_{yyy},g_{yyyy}\), supplied by the
complete moments, bound the remainder by
\(C(|y|^3/6+t|y|)+Dt^2/2\); the first two spatial derivative errors
are bounded by \(C(y^2/2+t)\) and \(C|y|+Dt\).
Substitution in \((g_y)^2-gg_{yy}\) and Gaussian integration with
variance \(v=\kappa t\) therefore prove
\begin{equation}\label{eq:Gaussian-event}
 A_{m-2}(\lambda_*-t,x_*,\kappa t)
 =a_m^2t(\kappa/2-1)+O_{a_m,C,D,\kappa}(t^{3/2}),\quad0<\kappa<2.
\end{equation}
The explicit remainder in the cited parent theorem pays the full
Gaussian exterior. At each fixed event the expression is negative
for sufficiently small \(t\). A uniform sign-prevention theorem would
still need a source-specific estimate at every relevant order and
shrinking variance, with event amplitude paid. The hierarchy is a
diagnostic, not such a theorem.

For \(n=0\), putting \(Z=\int w_vH_\lambda^2\) and
\(E_w=\int|(H_\lambda\sqrt{w_v})'|^2\) also gives
\(A_0=2E_w-Z/(2v)\), by integration by parts. Thus \(E_w\ge0\)
is weaker than the signed threshold \(E_w\ge Z/(4v)\). This is the
same distinction between regular positive energy and a signed
zero-sensitive observable encountered in Section~\ref{sec:geometry}.

\subsection{The concrete collision input maintained by OA}
The real-pair gap law \eqref{eq:q-exact} has a finite-interval consumer
in OA \cite[Section 4]{GNPCollisions}. On a predeclared box of times,
centers and virtual real node pairs, acquire the deflated factor by
divided differences with a uniform model floor and derivative errors.
These must imply \(|g|\le A\), \(|g_x|\le K\), and an upper bound
\(\beta\) for the averaged exterior curvature in \eqref{eq:q-exact}.
The floor is required throughout that box, before selecting the unknown
trajectory. Values only at the initial roots do not supply it.

Here \(c=(r_++r_-)/2\) is the pair center and \(D\) bounds its half-gap.
Suppose initially \(|c(0)-X|\le R_0<R\) and
\(0<Q_-\le q(0)\le Q_+<4D^2\). If
\begin{equation}\label{eq:pair-survival}
 R_0+2A\Delta<R,\qquad Q_+\ee^{4K\Delta}<4D^2,\qquad
 Q_--8\int_0^\Delta\ee^{4\beta t}\dd t>0,
\end{equation}
the specified pair stays real, distinct and simple for
\(0\le t\le\Delta\). Indeed \(|c'|\le2A\), \(q'\le4Kq\), and
\(q'\ge-8-4\beta q\) give confinement and the positive lower bound
\[
 q(t)\ge\ee^{-4\beta t}
       \left(Q_--8\int_0^t\ee^{4\beta s}\dd s\right).
\]
At a proposed stopping time the root slopes remain at least
\(\sqrt q\) times the factor floor, so the implicit function theorem
continues the pair. This proves the closed-interval assertion, including
the endpoint. A failed sufficient inequality leaves the actual pair
undecided. For the quadratic heat model the gap condition is sharp;
bounded exterior curvature by itself cannot prevent the terminal
\(-8\) closing rate.

The alternative contour interface extracts \(g=U'/U\) as
\((2\pi\ii)^{-1}\int_\Gamma(H'/H)(\zeta)/(\zeta-z)\dd\zeta\): the
cluster divisor's two residues cancel. Ratio errors on the boundary
then pay errors in \(g\) and \(g_x\) by the contour length and buffer
distance. This requires a complete cluster count and boundary margin,
but no lower internal gap. The same analytical tools therefore serve
Green acquisition and collision-force acquisition with different
consumers. A time-zero Green value alone does not provide the
space-time floor required in \eqref{eq:pair-survival}.
